\section*{Hoofdstuk \ref{ch:g12:glucose-and-diabetes} --- Bloedglucose en diabetes}

\begin{solution}{exo:g12:glucose-and-diabetes:1}
Ongeveer \qty{1}{g/L}, dus \qty{5.5}{mmol/L}. Onder \qty{0.5}{g/L} falen
de hersenen; boven \qty{1.8}{g/L} gaat er glucose in de urine (en jaren
boven \qty{1.3}{g/L} beschadigen de vaten).
\end{solution}

\begin{solution}{exo:g12:glucose-and-diabetes:2}
\omterm{prop:g12:glucose-and-diabetes:hormones}{Insuline}, uit de $\beta$-cellen, verlaagt de spiegel (opname en opslag
door spier, vet en lever); \omterm{prop:g12:glucose-and-diabetes:hormones}{glucagon}, uit de $\alpha$-cellen, verhoogt haar
(afgifte van glucose door de lever).
\end{solution}

\begin{solution}{exo:g12:glucose-and-diabetes:3}
De lever, als \omterm{def:g10:exercise-and-energy:glycogen}{glycogeen} (en zij kan ook nieuwe glucose maken uit
\omterm{def:g10:chemistry-of-life:families}{aminozuren} en lactaat).
\end{solution}

\begin{solution}{exo:g12:glucose-and-diabetes:4}
Type 1: de $\beta$-cellen zijn vernietigd en de \omterm{prop:g12:glucose-and-diabetes:hormones}{insuline} ontbreekt. Type
2: de \omterm{prop:g12:glucose-and-diabetes:hormones}{insuline} is er wel maar de weefsels reageren er slecht op, en de
$\beta$-cellen raken op den duur vermoeid.
\end{solution}

\begin{solution}{exo:g12:glucose-and-diabetes:5}
Gezond ongeveer \qty{0.95}{g/L}; verminderde tolerantie ongeveer
\qty{1.7}{g/L} (tussen 1.4 en 2); type 2 ongeveer \qty{2.7}{g/L} (boven 2).
\end{solution}

\begin{solution}{exo:g12:glucose-and-diabetes:6}
De lever geeft glucose uit haar \omterm{def:g10:exercise-and-energy:glycogen}{glycogeen} af met het tempo waarin de
hersenen haar wegnemen, onder de sturing van \omterm{prop:g12:glucose-and-diabetes:hormones}{glucagon}, zodat de voorraad
rond \qty{5}{g} blijft hoewel haar inhoud elk uur wordt ververst.
\end{solution}

\begin{solution}{exo:g12:glucose-and-diabetes:7}
De buis voert verteringsenzymen naar de darm; die afsluiten laat de
spiegel normaal, dus loopt het gevolg van de alvleesklier voor de glucose
niet via de vertering. Het hele orgaan weghalen haalt de \omterm{prop:g12:glucose-and-diabetes:hormones}{eilandjes} weg: de
\omterm{prop:g12:glucose-and-diabetes:hormones}{insuline} wordt in de alvleesklier gemaakt en bereikt haar doelwitten via
het bloed en niet via de buis.
\end{solution}

\begin{solution}{exo:g12:glucose-and-diabetes:8}
\omterm{prop:g12:glucose-and-diabetes:hormones}{Glucagon} werkt door de lever haar \omterm{def:g10:exercise-and-energy:glycogen}{glycogeen} te laten afbreken; na een
nacht vasten is de voorraad er en stijgt de spiegel; na drie dagen is het
\omterm{def:g10:exercise-and-energy:glycogen}{glycogeen} op en heeft \omterm{prop:g12:glucose-and-diabetes:hormones}{glucagon} niets meer om af te geven.
\end{solution}

\begin{solution}{exo:g12:glucose-and-diabetes:9}
Boven ongeveer \qty{1.8}{g/L} kunnen de nieren niet alle glucose die zij
filteren weer opnemen; die gaat in de urine en trekt door osmose water
mee. Het verloren water moet worden aangevuld: dorst.
\end{solution}

\begin{solution}{exo:g12:glucose-and-diabetes:10}
De \omterm{prop:g12:glucose-and-diabetes:hormones}{insuline} drijft glucose de spieren en de vetcellen in en stopt de
afgifte door de lever, maar er komt geen glucose binnen: de spiegel zakt
binnen een uur of twee onder \qty{0.6}{g/L}, met zweten, beven en
verwarring; meteen suiker nemen verhelpt dat.
\end{solution}

\begin{solution}{exo:g12:glucose-and-diabetes:11}
Een werkende spier neemt glucose uit het bloed op langs een weg die geen
\omterm{prop:g12:glucose-and-diabetes:hormones}{insuline} nodig heeft, en een getrainde spier wordt tussen de sessies weer
gevoeliger voor \omterm{prop:g12:glucose-and-diabetes:hormones}{insuline}; de spiegel daalt tijdens en na de inspanning, en
de resistentie van de weefsels, de oorzaak van de ziekte, wordt kleiner.
\end{solution}

\begin{solution}{exo:g12:glucose-and-diabetes:12}
Geen \omterm{prop:g12:glucose-and-diabetes:hormones}{insuline}: type 1, de $\beta$-cellen van de \omterm{prop:g12:glucose-and-diabetes:hormones}{eilandjes} zijn vernietigd
en scheiden niets af. Drie keer de normale \omterm{prop:g12:glucose-and-diabetes:hormones}{insuline}: type 2, de \omterm{prop:g12:glucose-and-diabetes:hormones}{eilandjes}
scheiden voluit af tegen weefsels die niet meer reageren.
\end{solution}

\begin{solution}{exo:g12:glucose-and-diabetes:13}
Gezond: de spiegel stijgt na 30--60 minuten tot ongeveer \qty{1.3}{g/L} en
is na 2 uur terug; de \omterm{prop:g12:glucose-and-diabetes:hormones}{insuline} stijgt en daalt met haar mee. Onbehandeld
type 1: de spiegel klimt boven \qty{2}{g/L} en blijft daar; de \omterm{prop:g12:glucose-and-diabetes:hormones}{insuline}
blijft op nul --- niets haalt de glucose weg. Type 2: de spiegel stijgt
hoog en daalt traag; de \omterm{prop:g12:glucose-and-diabetes:hormones}{insuline} stijgt hoger dan normaal en blijft hoog
--- veel \omterm{def:g11:hormones-and-reproduction:hormone}{hormoon}, weinig gevolg.
\end{solution}

\begin{solution}{exo:g12:glucose-and-diabetes:14}
$5 + 80 = \qty{85}{g}$ in \qty{5}{L}: \qty{17}{g/L}. Dat gebeurt nooit
omdat de opname een uur of twee duurt, en terwijl de glucose binnenkomt
stuurt de \omterm{prop:g12:glucose-and-diabetes:hormones}{insuline} haar naar het leverglycogeen, het spierglycogeen en het
vet, terwijl de hersenen en de weefsels er wat van verbranden: de voorraad
bevat nooit meer dan een paar gram boven de normale waarde.
\end{solution}

\begin{solution}{exo:g12:glucose-and-diabetes:15}
Beide: een voeler (\omterm{def:g10:cells-common-unit:cell}{cellen} van de \omterm{prop:g12:glucose-and-diabetes:hormones}{eilandjes}; hypothalamus en hypofyse), een
streefwaarde, uitvoerders, en negatieve \omterm{def:g11:hormones-and-reproduction:hormone}{terugkoppeling} die het signaal
wegneemt. De \omterm{def:g12:glucose-and-diabetes:glycaemia}{bloedglucose} heeft twee tegengestelde \omterm{def:g11:hormones-and-reproduction:hormone}{hormonen} die op
dezelfde organen werken, zodat de grootheid actief zowel omhoog als omlaag
kan worden geduwd; de lus van het testosteron heeft één \omterm{def:g11:hormones-and-reproduction:hormone}{hormoon} en stuurt
vooral bij met meer of minder daarvan. Twee tegengestelde uitvoerders
geven een snellere, strakkere sturing van een grootheid die in beide
richtingen snel verandert.
\end{solution}

\begin{solution}{pb:g12:glucose-and-diabetes:1}
\textbf{1.} \qty{5}{g}; $\qty{1}{g/L} \div \qty{180}{g/mol} = \qty{5.55}{mmol/L}$.

\textbf{2.} Eén uur.

\textbf{3.} $100/5 = 20$ uur. Daarna maakt de lever nieuwe glucose uit
\omterm{def:g10:chemistry-of-life:families}{aminozuren} (en lactaat), ten koste van de \omterm{def:g10:chemistry-of-life:families}{eiwitten} van het lichaam.

\textbf{4.} Spiercellen missen het \omterm{def:g11:enzymes-and-phenotype:enzyme}{enzym} dat vrije glucose uit \omterm{def:g10:exercise-and-energy:glycogen}{glycogeen}
aan het bloed afgeeft; hun \omterm{def:g10:exercise-and-energy:glycogen}{glycogeen} wordt ter plekke verbrand, voor hun
eigen samentrekking.

\textbf{5.} $5 + 60 = \qty{65}{g}$ in \qty{5}{L}: \qty{13}{g/L}. De
\omterm{prop:g12:glucose-and-diabetes:hormones}{insuline} stuurt de glucose naar het lever- en spierglycogeen en naar het
vet, en de weefsels verbranden er wat van; de spiegel piekt rond \qty{1.4}{g/L}.

\textbf{6.} Nuchter \qty{1.15}{g/L} (onder 1.26) maar \qty{1.85}{g/L} op
120 minuten (tussen 1.4 en 2): verminderde tolerantie, op de drempel van
\omterm{def:g12:glucose-and-diabetes:diabetes}{diabetes}.

\textbf{7.} Nee: de \omterm{prop:g12:glucose-and-diabetes:hormones}{insuline} zit op twee keer een gezonde piek. De reactie
van de weefsels erop faalt --- insulineresistentie.

\textbf{8.} Type 2 (of het waarschuwende stadium). Verwacht: volwassen,
overgewicht, weinig beweging, vaak met suikerzieke familieleden.

\textbf{9.} Een spiegel die boven \qty{1.5}{g/L} begint, boven \num{2.5}
stijgt en niet zakt; de \omterm{prop:g12:glucose-and-diabetes:hormones}{insuline} blijft de hele tijd bij nul.

\textbf{10.} Afvallen en geregeld bewegen: zij herstellen de gevoeligheid
van de weefsels voor \omterm{prop:g12:glucose-and-diabetes:hormones}{insuline} en laten werkende spieren glucose opnemen
zonder haar.

\textbf{11.} $90/10 = 9$ eenheden.

\textbf{12.} De dosis verwerkt \qty{90}{g} maar er komt maar \qty{45}{g}
binnen: de overtollige \omterm{prop:g12:glucose-and-diabetes:hormones}{insuline} duwt de spiegel omlaag met het equivalent
van \qty{45}{g} in de voorraad en de opslag --- ruim onder \qty{0.6}{g/L}:
zweten, beven, verwarring; meteen suiker.

\textbf{13.} $(5 + 90)/5 = \qty{19}{g/L}$ als er niets werd weggehaald; in
werkelijkheid enkele \unit{g/L}. Boven \qty{1.8}{g/L} laten de nieren
glucose in de urine, met water mee: veel urine en dorst.

\textbf{14.} Tussen de maaltijden zou de lever anders ongeremd glucose
afgeven (\omterm{prop:g12:glucose-and-diabetes:hormones}{glucagon} werkt, niets houdt het tegen) en zouden de weefsels vet
verbranden met zure producten: er is een achtergrond van \omterm{prop:g12:glucose-and-diabetes:hormones}{insuline} nodig
voor de basisbalans, niet alleen voor de maaltijden.

\textbf{15.} De meter meet de \omterm{def:g12:glucose-and-diabetes:glycaemia}{bloedglucose} zoals de $\beta$-cel dat deed;
de pen levert de \omterm{prop:g12:glucose-and-diabetes:hormones}{insuline} die zij zou hebben afgescheiden. Zij kunnen de
voortdurende bijstelling van minuut tot minuut niet vervangen: de patiënt
meet een paar keer per dag en schat, terwijl de \omterm{def:g10:cells-common-unit:cell}{cel} onophoudelijk mat en
precies reageerde.

\textbf{16.} Type 2, waarvan de $\beta$-cellen er zijn en kunnen worden
aangezet. Bij type 1 zijn er geen $\beta$-cellen om aan te zetten.

\textbf{17.} De reactie van de weefsels op \omterm{prop:g12:glucose-and-diabetes:hormones}{insuline} --- de uitvoerderskant
van de lus: met minder vet en actieve spieren geeft een normaal
insulinesignaal weer een normale opname, en de $\beta$-cellen krijgen lucht.

\textbf{18.} Elk wordt afgescheiden als antwoord op de tegengestelde
afwijking van dezelfde grootheid: een hoge spiegel roept \omterm{prop:g12:glucose-and-diabetes:hormones}{insuline} en legt
\omterm{prop:g12:glucose-and-diabetes:hormones}{glucagon} stil, een lage andersom. Bij \qty{1}{g/L} scheidt het \omterm{prop:g12:glucose-and-diabetes:hormones}{eilandje} ze
allebei in een laag, evenwichtig tempo af --- zijn rusttoestand.

\textbf{19.} Een lage spiegel doodt de hersenen binnen enkele minuten, een
hoge beschadigt traag over jaren: het acute gevaar heeft verscheidene
onafhankelijke wachters, het trage één --- en daarom is het uitvallen van
de \omterm{prop:g12:glucose-and-diabetes:hormones}{insuline} alleen al genoeg om een ziekte te veroorzaken, en daarom is
\omterm{def:g12:glucose-and-diabetes:diabetes}{diabetes} veelvoorkomend.

\textbf{20.} Ongeveer \qty{5}{g} glucose in het bloed; het \omterm{def:g10:exercise-and-energy:glycogen}{glycogeen} van
de lever overbrugt ongeveer 20 uur; de patiënt van deel II heeft een
verminderde tolerantie op de rand van \omterm{def:g12:glucose-and-diabetes:diabetes}{diabetes} type 2; de dosis van deel
III is 9 eenheden.
\end{solution}
